\chapter{Methodology}

This chapter describes the experimental design, workload selection, measurement approach, and statistical methodology used to compare multi-threaded and distributed parallelism models for matrix operations.

\section{Research Method}

A controlled experiment isolates the parallelism model as the independent variable while holding constant all other factors: workload type, total compute capacity, matrix dimensions, numerical kernel, and hardware platform.

\subsection{Independent Variables}

\textbf{Parallelism Model} (3 levels):
\begin{itemize}
    \item \textbf{Threading}: Python \texttt{threading} module with shared memory
    \item \textbf{Multiprocessing}: Python \texttt{multiprocessing} module with separate address spaces
    \item \textbf{Distributed}: Dask distributed scheduler with \texttt{LocalCluster} on one host (local proxy for scale-out; \textbf{not} multi-instance AWS EC2)
\end{itemize}

\textbf{Matrix Size} (factor with 5 levels): 200×200, 500×500, 1000×1000, 1500×1500, 2000×2000

\textbf{Worker Count} (factor with 4 levels): 1, 2, 4, 8 workers/threads/processes

\subsection{Dependent Variables}

\begin{enumerate}
    \item \textbf{Completion Time}: Wall-clock time from operation start to finish (seconds). Includes initialization and data distribution overhead as these costs are incurred by practitioners.
    
    \item \textbf{Peak Memory}: Maximum resident set size (RSS) across all processes during execution (MB). Sampled at 50ms intervals.
    
    \item \textbf{CPU Utilization}: Mean percentage of available CPU capacity used during execution. Sampled at 100ms intervals. A value of 400\% on a 4-core system indicates full utilization.
\end{enumerate}

\subsection{Controlled Variables}

\begin{itemize}
    \item \textbf{Total vCPU Count}: When comparing parallelism models, the total number of cores is held constant. For example, 4-worker distributed is compared to 4-thread multi-threaded, not to 2-thread.
    
    \item \textbf{Numerical Kernel}: NumPy 1.26.4 with OpenBLAS backend. Internal threading is controlled (\texttt{OMP\_NUM\_THREADS=1} for threading model to isolate our parallelism).
    
    \item \textbf{Matrix Generation}: Fixed random seed (42) ensures identical input data across runs.
    
    \item \textbf{Hardware Platform}: All experiments execute on the same physical machine to eliminate hardware variability.
    
    \item \textbf{Operating System State}: Background processes minimized; no other intensive workloads running.
\end{itemize}

\section{Workload Selection}

\subsection{Matrix Multiplication}

Dense matrix multiplication $C = A \times B$ for $n \times n$ matrices has $O(n^3)$ computational complexity and $O(n^2)$ memory requirement. It is:

\begin{itemize}
    \item Compute-intensive with predictable cost
    \item Well-supported by mature parallel libraries
    \item Emblematic of scientific computing workloads
    \item Used as a standard parallel computing benchmark
\end{itemize}

The operation parallelizes naturally by partitioning rows of $A$ across workers, each computing its assigned rows of $C$ via:
\[
C[i,:] = A[i,:] \times B
\]

\subsection{LU Factorization}

LU decomposition factors a matrix $A = LU$ where $L$ is lower triangular and $U$ is upper triangular, using Gaussian elimination. Complexity is $O(n^3)$ for dense matrices.

This workload is chosen to match Sabir and Alebrahim (2025) \cite{sabir2025coarse}, enabling direct comparison with their multi-threaded baseline. LU factorization is:

\begin{itemize}
    \item Sequential in nature (elimination progresses column-by-column)
    \item More challenging to parallelize effectively
    \item Sensitive to synchronization overhead
    \item Critical for solving linear systems
\end{itemize}

For numerical stability, all test matrices are positive definite: $A_{pd} = MM^T + nI$ where $M$ is random and $n$ is the matrix dimension.

\section{Experimental Design}

\subsection{Configuration Matrix}

The full experimental design comprises:

\begin{itemize}
    \item 5 matrix sizes × 4 worker counts = 20 size/worker combinations
    \item 3 parallelism models (threading, multiprocessing, distributed)
    \item 2 operations (matrix multiplication, LU factorization)
    \item Note: Distributed requires ≥2 workers, so 1-worker distributed configurations are omitted
\end{itemize}

Committed unique configurations in \texttt{summary\_statistics.json}: \textbf{90} (sizes 200--2000; workers 1--8; distributed omits 1-worker cells).

\subsection{Replication and Statistical Power}

\textbf{Planned} replication was 5 iterations per configuration (pilot CV and nominal 80\% power at $\alpha=0.05$). \textbf{Committed} results use \texttt{n\_iterations: 3} for every row in \texttt{summary\_statistics.json} (90 configs × 3 = 270 runs). Report numbers must follow the committed $n{=}3$ evidence, not the planned $n{=}5$ text.

\subsection{Measurement Protocol}

For each configuration:

\begin{enumerate}
    \item Generate test matrices with fixed seed
    \item Start memory and CPU monitors (background threads sampling at intervals)
    \item Record start time (\texttt{time.perf\_counter})
    \item Execute matrix operation
    \item Record end time
    \item Stop monitors and collect peak/mean values
    \item Save result to JSON/CSV
\end{enumerate}

The measurement includes all overhead (cluster setup, data serialization, worker coordination) as these are costs practitioners must pay.

\section{Statistical Analysis}

\subsection{What is committed vs planned}

Committed artefacts report \textbf{descriptive} statistics only (per-config mean and standard deviation over $n{=}3$ iterations in \texttt{summary\_statistics.json}). The following inferential plan remains \textbf{methodology intent}; corresponding Shapiro / $t$-test / Mann--Whitney / Holm code and $p$-value artefacts are \textbf{not} present under \texttt{distributed-matrix-scaling/} and are therefore not claimed as executed evidence.

\subsection{Null and Alternative Hypotheses (planned)}

For each matrix size and worker count:

\begin{itemize}
    \item $H_0$: No difference in completion time between parallelism models
    \item $H_1$: Significant difference exists
\end{itemize}

Similarly for memory consumption and CPU utilization.

\subsection{Test Selection (planned, not coded)}

\begin{enumerate}
    \item \textbf{Normality Test}: Shapiro-Wilk test on residuals. If $p > 0.05$, proceed with parametric tests.
    
    \item \textbf{Parametric}: Independent-samples t-test for pairwise comparisons, with Cohen's d effect size.
    
    \item \textbf{Non-parametric}: Mann-Whitney U test if normality assumption is violated, with rank-biserial correlation as effect size.
\end{enumerate}

\subsection{Multiple Comparisons Correction (planned, not coded)}

Holm-Bonferroni correction would control family-wise error at $\alpha=0.05$ if inferential tests were implemented.

\subsection{Crossover Identification}

A crossover point is defined as the smallest matrix size where:
\begin{enumerate}
    \item Distributed model shows faster mean completion time than multi-threaded at matched workers
    \item This advantage persists for all larger matrix sizes tested
\end{enumerate}

In the committed local JSON, this condition is \textbf{never} met for matmul through 2000×2000.

\section{Execution Environment}

\subsection{Hardware and Software}

\begin{itemize}
    \item \textbf{Hardware}: [Actual system specs documented in results metadata]
    \item \textbf{OS}: Ubuntu 24.04 LTS
    \item \textbf{Python}: 3.12.3
    \item \textbf{NumPy}: 1.26.4 with OpenBLAS
    \item \textbf{Dask}: 2026.8.0
    \item \textbf{Distributed}: 2026.8.0
\end{itemize}

\subsection{Local vs. Cloud Execution (CA2 requirement vs practice)}

\textbf{CA2 requirement:} matched aggregate vCPU comparison on AWS EC2 (multi-threaded on a single instance vs distributed across instances), with infrastructure provisioned and destroyed via IaC.

\textbf{Practice in this deliverable:} all experiments run on a \textbf{local} development machine with Dask \texttt{LocalCluster}. Local matched cores on one host are \textbf{not} equivalent to matched aggregate vCPUs across EC2 instances. No live EC2 / CloudWatch evidence is claimed.

Rationale for the local proxy:

\begin{itemize}
    \item \textbf{Reproducibility}: No AWS account or credentials required; anyone can run the benchmarks
    \item \textbf{Cost}: Zero infrastructure cost during development and validation
    \item \textbf{Honesty}: Local absolute times must not be presented as cloud results
\end{itemize}

The local Dask cluster (multiple workers on same machine) is only a proxy for scale-out:

\begin{itemize}
    \item Distributed overhead is likely \textit{underestimated} relative to multi-node AWS (no real network fabric)
    \item Any future crossover on EC2 must be measured; it is not evidenced here
    \item Remote \texttt{--scheduler-address} CLI wiring is not present in committed \texttt{main.py}
\end{itemize}

\section{Ethical Considerations}

This study:
\begin{itemize}
    \item Involves no human participants
    \item Processes no personal or sensitive data
    \item Uses only synthetic matrices generated from known seeds
    \item Consumes compute resources within acceptable-use bounds (local execution, limited duration)
    \item Reports measurements honestly without selective reporting or p-hacking
\end{itemize}

All source code, raw data, and analysis scripts are provided for transparency and reproducibility.

\section{Threats to Validity}

\subsection{Internal Validity}

\begin{itemize}
    \item \textbf{Background Processes}: Mitigated by closing non-essential applications and running experiments during idle periods
    \item \textbf{Thermal Throttling}: Possible for sustained workloads; mitigated by monitoring CPU frequency and spacing runs
    \item \textbf{OS Scheduling}: Cannot be fully controlled; mitigated by multiple iterations and statistical analysis
\end{itemize}

\subsection{External Validity}

\begin{itemize}
    \item \textbf{Hardware Specificity}: Results are platform-dependent; trends generalize but absolute times do not
    \item \textbf{Matrix Operations}: Findings apply to dense matrix workloads; sparse matrices or other algorithms may differ
    \item \textbf{Local Execution}: True cloud network latency is not captured; distributed overhead is likely understated
\end{itemize}

\subsection{Construct Validity}

\begin{itemize}
    \item \textbf{Measurement Overhead}: Memory/CPU sampling introduces minimal overhead (<1\% estimated)
    \item \textbf{Cold Start Effects}: First iteration may include JIT compilation; all configurations affected equally
\end{itemize}

\section{Summary}

The experimental methodology controls for hardware, software, and workload factors while systematically varying parallelism model, matrix size, and worker count. Committed evidence is descriptive ($n{=}3$ means/std in JSON). Local execution trades true cloud network characteristics for reproducibility and zero cost; the CA2 AWS EC2 matched-vCPU campaign remains outstanding.
